\songtitle{Petrovna // Петровна}

\tag*{2026}\tag{folk-metal}\tag{thomas-entourage}\tag{kadri-kutsnisova}\tag{estonian-lyrics}\tag{russian-lyrics}

\begin{notes}
\end{notes}

\begin{style}
female vocals, female vocals! important: clean folk vocals in verses, Folk/pagan metal, Mid-paced, processional — weight over speed, Foundation: heavy, droning down-tuned guitars and bass, slow rhythmic pulse, no blast beats, Color layer: kantele or gusli motif in intro and bridge, threading between the heavy sections like something older than the arrangement, Drumming: deliberate, tribal-adjacent in verses, fuller and driving in choruses, Vocal: clean, restrained folk delivery in Estonian verses — precise, slightly detached; fuller, rawer, more resonant in Russian choruses, as if a different register of the same person, Bridge: vocal stripped back to near-spoken, kantele prominent, guitars recede, Final chorus: fullest delivery, guitars return at full weight, Outro: drone and kantele only, fading
\end{style}

\begin{lyrics}
\songpart{Intro}
\annotation{kantele over slow drone — 16 bars, no vocal}

\songpart{Verse 1}
\annotation{clean folk vocal, guitars enter quietly beneath}
\annotation{Estonian}\\
Sündinud seal, kus vanalinn merele vaatab\\
isa veri idast, ema käed lõunast\\
kaks keelt enne kooli, kaks viisi külma tunda
õppisin kõik suu sees hoidma

Tallinn jäi, teised lahkusid\\
tegin kodu sellest, mis jäi

\songpart{Pre-Chorus}
\annotation{guitars build, drums enter, vocal firms up}
\annotation{Estonian}\\
Loen neid, kes tulevad siia eksinult\\
tean märke enne, kui nad räägivad\\
ehitasin oma seinad sellest, mida leidsin\\
ja see, mida leidsin — mul pole vaja rohkem

\songpart{Chorus}
\annotation{full band, driving pulse, vocal raw and resonant}
\annotation{Russian}
Мне не нужен один
я держу нескольких близко
не из страха остаться
а так я устроена

Тепло без якоря
близость без границ
я знаю себя —
мне этого хватит

\songpart{Verse 2}
\annotation{clean folk vocal returns, guitars drop back to verse weight}
\annotation{Estonian}\\
Valisin need, kes lagunevad\\
võõrastes linnades, võõrastes keeltes\\
loen toimikut enne, kui räägin\\
tean, mida nad otsivad, enne kui nad teavad

Nahktagi ja rasused juuksed\\
rebitud kaelaava — see on valik, mitte hoolimatus\\
mitte peitmine — ütlemine: vaata, vaata uuesti\\
ma näitan sulle, mida tahan, et sa näeksid

\songpart{Pre-Chorus 2}
\annotation{guitars build again, drums fuller}
\annotation{Estonian}

Seinad pole haavad —\\
need on arhitektuur\\
ehitasin need teadlikult\\
tuba endale\\
toad teistele\\
uksed avanevad minu tahtel

\songpart{Chorus}
\annotation{full band}
\annotation{Russian}
Мне не нужен один
я держу нескольких близко
не из страха остаться
а так я устроена

Тепло без якоря
близость без границ
я знаю себя —
мне этого хватит

\songpart{Bridge}
\annotation{guitars pull back to drone only; kantele returns prominent; vocal near-spoken, unhurried; lines intercut between languages with space between them}
\annotation{Russian}  Но иногда кто-то смотрит иначе
\annotation{Estonian} ja ma ei ole valmis selleks
\annotation{Russian}  не сквозь меня — на меня
\annotation{Estonian} see on midagi muud
\annotation{Russian}  и я позволяю себе это заметить
\annotation{Estonian} ainult märgata
\annotation{Russian}  не больше

\songpart{Final Chorus}
\annotation{guitars return at full weight; heaviest, most driven moment of the song; vocal at fullest delivery throughout}
\annotation{Russian}
Мне не нужен один
я держу нескольких близко
не из страха остаться
а так я устроена

Тепло без якоря
близость без различий
я знаю себя —
\annotation{Estonian, half-sung over final bars, guitars sustaining}\\
see on minu oma

\songpart{Outro}
\annotation{drone and kantele only; guitars fade on final chord; kantele carries the melodic thread alone until silence}
\end{lyrics}

